针对通用神经网络处理器中拓扑依赖、有限连续缓存与多执行单元并行的耦合，建立以生命周期收缩、驻留段分配和固定驻留图重排为核心的调度方法。在六个题面计算图上，分别求解最小驻留、低额外搬运和搬运预算下的执行时间优化。下述六图结果依次对应卷积小图、大图，注意力小图、大图及矩阵乘小图、大图。

对于问题一，以缓冲区活跃区间建立峰值目标，证明在申请为根、释放为叶的条件下，固定操作顺序的生命期收缩不增加峰值，再比较编号、窗口贪心与深度优先等八种顺序。卷积小图峰值由 7560 降至 6984，其余五图分别为 14048、4884、7972、9728、35328。通过九节点小例的 126 个合法拓扑序验证收缩过程，并用分类型驻留轨迹解释局部贪心与完整峰值的差异。

对于问题二，用四状态描述逻辑生命与物理驻留，结合连续区间最佳适配及未来使用距离选择换出对象；受保护对象形成碎片时，以显式换出和重装恢复空间，并计入全部代价。六图额外搬运量依次为 44766、75546、3620、32852、34944、460800，卷积小图相对编号基准降低 48.90\%。同顺序的换出准则对照表明，考虑单次搬运代价并不普遍优于单纯未来距离。

对于问题三，在搬运量零增量约束下比较地址轮换，并在含地址复用依赖的固定驻留图上实施列表重排。证明该重排保留地址互斥及换出集合，最终时间分别为 465892、1044075、37744、150499、173274、1749773 周期，相对问题二降低 1.30\% 至 33.38\%。注意力大图的向量单元忙碌比例由 46.46\% 提高至 69.74\%。进一步区分原始工作下界与固定驻留下界，六图相对后者的差距为 1.53\% 至 24.95\%，仅用于评价保留当前驻留关系后的改进空间。

所有方案经原始依赖、连续地址、使用时驻留和执行单元互斥重算。分类型搬运、重复换出及局部流水对照表明，缓存压力集中位置与地址链结构共同影响三问的优化收益。结果适用于六个题面计算图和本文采用的驻留段复用模型。
